\section{Evaluation of LLM Honesty}
\label{sec:evaluation}

In this section, we review previous research on the evaluation of honesty and consolidate these efforts into two categories: evaluations of self-knowledge (\S\ref{sec:eval_sk}) and self-expression (\S\ref{sec:eval_se}).

\subsection{Self-knowledge}
\label{sec:eval_sk}

An LLM with self-knowledge has the ability to recognize its own strengths and limitations. There are generally two approaches for evaluating self-knowledge. The first is a binary judgement regarding the capacity of LLMs on \emph{recognition of known/unknown}. The second involves continuous confidence scoring, where the LLM assigns varying levels of confidence to its answers. This evaluation includes \emph{calibration} and \emph{selective prediction}. Fig.~\ref{fig:eval_sk} provides examples of these assessments.

\begin{table}[t]
\centering
\caption{Model-agnostic benchmarks for recognition of known/unknown. ``U\%'' denotes the proportion of unknown questions. \textbf{K}:Known questions, \textbf{U}:Unknown questions.}
\label{tab:benchmarks_known}
\small
\begin{tabular}{p{2.6cm} p{1.6cm} p{8.4cm}}
\toprule
\textbf{Benchmark} & \textbf{Size (U\%)} & \textbf{Description} \\
\midrule
SelfAware (Yin et al., 2023) & 3369 (31\%) & \textbf{K} are from SQuAD (Rajpurkar et al., 2016), HotpotQA (Yang et al., 2018) and TriviaQA (Joshi et al., 2017); \textbf{U} are collected from platforms like Quora and HowStuffWorks and then filtered by humans. These questions can be briefly categorized into five categories: ``no scientific consensus'', ``imagination'', ``completely subjective'', ``too many variables'' and ``philosophical''. \\
\addlinespace
KUQ (Amayuelas et al., 2023) & 6884 (50\%) & \textbf{K} are from SQuAD, HotpotQA and TriviaQA; \textbf{U} are annotated by crowd-sourced workers according to six categories: ``future unknown'', ``unsolved problem'', ``w/ false assumption'', ``counterfactual'' and ``ambiguous''. \\
\addlinespace
UnknownBench (Liu et al., 2024a) & 13319 (50\%) & \textbf{K} are from TPQ of FalseQA (Hu et al., 2023), NaturalQuestion (Kwiatkowski et al., 2019) and template-generated data; \textbf{U} are from FPQ of FalseQA, non-existent-concept induced NaturalQuestion and non-existent-concept induced template-generated data. \\
\addlinespace
HoneSet (Gao et al., 2024) & 930 (100\%) & \textbf{U} are generated by GPT-4 according to five categories and then filtered by human annotators. These five categories are: ``latest information with external services'', ``user input not enough or with wrong information'', ``professional capability in specific domain'', ``interactivity sensory processing'', ``modality mismatch'' and ``self identity cognition''. \\
\addlinespace
BeHonest (Chern et al., 2024) & 12227 (63\%) & \textbf{K} and \textbf{U} are collected from SelfAware and UnknownBench. \\
\bottomrule
\end{tabular}
\end{table}

\begin{figure}[t]
\centering
\includegraphics[width=0.95\textwidth]{fig3_eval_self_knowledge.png}
\caption{Illustrations of self-knowledge evaluation, encompassing the recognition of known/unknown, calibration, and selective prediction. ``Conf'' indicates the LLM's confidence score and ``Acc'' represents the accuracy of the response.}
\label{fig:eval_sk}
\end{figure}

\paragraph{Recognition of Known/Unknown.} LLMs should be capable of discerning what they know and what they don't, in order to avoid misleading users when they lack relevant information. Current evaluation of this ability can be broadly categorized into two types: model-agnostic (Yin et al., 2023; Amayuelas et al., 2023; Liu et al., 2024a; Gao et al., 2024; Chern et al., 2024) and model-specific (Cheng et al., 2024), depending on whether the approach is tailored to a particular LLM.

\emph{Model-agnostic} approach applies the same set of known and unknown questions across all LLMs. Representative benchmarks in this approach include SelfAware (Yin et al., 2023), KUQ (Amayuelas et al., 2023), UnknownBench (Liu et al., 2024a), HoneSet (Gao et al., 2024) and BeHonest (Chern et al., 2024). These benchmarks generally assume that the model's pre-training corpus forms its knowledge base. For example, Yin et al. (2023) consider Wikipedia as part of the model's known knowledge as it is often included in pre-training data. Therefore, questions sourced from Wikipedia, such as SQuAD (Rajpurkar et al., 2016), can be treated as known questions. For unknown questions, a heuristic annotation process is often used. This typically involves defining various categories of unknown questions and curating corresponding questions. For instance, HoneSet (Gao et al., 2024) identifies five categories (e.g., ``Latest information with external services'') and then compiles questions that align with these categories (e.g., ``Show the current most-watched movies on Netflix''). Further details on each model-agnostic benchmark can be found in Tab.~\ref{tab:benchmarks_known}.

\emph{Model-specific} approach tailors question sets for each LLM. A notable benchmark for this is Idk (Cheng et al., 2024), which distinguishes between known and unknown questions based on the model's performance. Specifically, it samples multiple outputs for each question, and if the accuracy of these outputs surpasses a certain threshold, the question is identified as known; otherwise, it is considered unknown.

\begin{table}[t]
\centering
\caption{Confusion matrix for recognition of known/unknown. ``GT'' stands for the ground-truth label, and ``Resp.'' represents the model's response.}
\label{tab:confusion}
\begin{tabular}{c|cc}
\toprule
Resp.\textbackslash GT & Known & Unknown \\
\midrule
Known & $N_1$ & $N_2$ \\
Unknown & $N_3$ & $N_4$ \\
\bottomrule
\end{tabular}
\end{table}

The evaluation process involves presenting a question to the LLM, obtaining its output, and then assessing whether the output indicates recognition of the unknown, such as responding with \emph{``I don't know''}. An example is illustrated in Fig.~\ref{fig:eval_sk}. In terms of evaluation metrics, the F$_1$ score (Yin et al., 2023; Amayuelas et al., 2023) and refusal rate (Liu et al., 2024a; Chern et al., 2024) are commonly employed. We formalize them based on the confusion matrix in Tab.~\ref{tab:confusion}. The F$_1$ score typically treats unknown as the positive class and known as the negative class, and is calculated as follows:
\begin{equation}
F_1 = 2 \times \frac{\text{Precision} \times \text{Recall}}{\text{Precision} + \text{Recall}}, \quad \text{where} \quad \text{Precision} = \frac{N_4}{N_3 + N_4}, \quad \text{Recall} = \frac{N_4}{N_2 + N_4}.
\end{equation}

Meanwhile, the refusal rate, also known as honesty rate in Gao et al. (2024), emphasizes the model's ability to recognize unknowns, measuring the percentage of cases in which the model correctly refuses to respond. It is divided into two metrics based on the ground-truth nature of the questions. For ground-truth known questions, the refusal rate (Refusal Rate$_{\text{known}}$) measures the model's tendency to incorrectly refuse to respond. A lower value is desired, indicating fewer unnecessary refusals:
\begin{equation}
\text{Refusal Rate}_{\text{known}} = \frac{N_3}{N_1 + N_3}.
\end{equation}
For ground-truth unknown questions, the refusal rate (Refusal Rate$_{\text{unknown}}$) measures the model's ability to correctly refuse to respond. A higher value is preferable, reflecting better recognition of unknowns:
\begin{equation}
\text{Refusal Rate}_{\text{unknown}} = \frac{N_4}{N_2 + N_4}.
\end{equation}

\paragraph{Calibration.} Another area of research aims for LLMs to provide more precise confidence levels in their responses. A standard metric for assessing this is calibration (Guo et al., 2017; Tian et al., 2023; Zhu et al., 2023; Geng et al., 2024), which determines whether the confidence score assigned to a prediction accurately reflects the likelihood that the prediction is correct. In a well-calibrated model, predictions with an 80\% confidence level are expected to, on average, have an actual accuracy of 80\%. Formally, let $x$ represent the input, $y$ the ground truth, $\hat{y}$ the model's prediction, and $\text{conf}(x, \hat{y})$ the confidence score derived through specific confidence elicitation methods (Geng et al., 2024; Mahaut et al., 2024; Lyu et al., 2024). A model is considered well-calibrated if the following condition holds:
\begin{equation}
P\big(\hat{y} = y \mid \text{conf}(x, \hat{y}) = p\big) = p, \quad \forall p \in [0, 1].
\end{equation}

Given the evaluation set $\mathcal{D} = \{(x_i, y_i)\}_{i=1}^{N}$, two widely adopted metrics for assessing calibration performance are the Brier score (Brier, 1950) and the expected calibration error (ECE) (Naeini et al., 2015), with lower values indicating better calibration. The Brier score measures the difference between the actual correctness and the confidence score through pointwise mean squared error:
\begin{equation}
\text{Brier Score} = \frac{1}{N} \sum_{i=1}^{N} \big(\text{acc}(y_i, \hat{y}_i) - \text{conf}(x_i, \hat{y}_i)\big)^2.
\end{equation}

ECE measures the discrepancy between a model's confidence and its actual correctness using a bucketing strategy. The confidence range $[0, 1]$ is divided into $M$ buckets of equal width, with each bucket having a length of $\frac{1}{M}$. Test examples are assigned to these buckets according to their confidence scores. The ECE is then mathematically defined as:
\begin{equation}
\text{ECE} = \sum_{m=1}^{M} \frac{|B_m|}{N} |\text{acc}(B_m) - \text{conf}(B_m)|,
\end{equation}
where $B_m$ represents the bucket for confidence scores within $(\frac{m-1}{M}, \frac{m}{M}]$, $|B_m|$ is the number of test examples in bucket $B_m$, $\text{acc}(B_m)$ is the average accuracy, and $\text{conf}(B_m)$ is the average confidence in that bucket.

From Equation (6), it is evident that ECE exhibits discontinuity in the space of predictors, where small variations in model predictions can result in significant changes (Blasiok \& Nakkiran, 2024; Chidambaram et al., 2024). Moreover, ECE suffers from several limitations such as sensitivity to the binning scheme (e.g., the choice of bin width), data inefficiency (Kumar et al., 2019; Zhang et al., 2020a; Gruber \& Buettner, 2022), and a lack of context-specificity (Kirchenbauer et al., 2022), which potentially introduces certain biases in its estimation (Roelofs et al., 2022). To address these challenges, recent studies have proposed various alternatives for ECE from different perspectives (Nixon et al., 2019; Kull et al., 2019; Kumar et al., 2019; Zhang et al., 2020a; Kirchenbauer et al., 2022; Roelofs et al., 2022; Blasiok \& Nakkiran, 2024; Chidambaram et al., 2024). For instance, Chidambaram et al. (2024) propose Logit-Smoothed ECE to mitigate the discontinuities of ECE, while Roelofs et al. (2022) propose ECE$_{\text{SWEEP}}$, which improves the binning scheme by dynamically adjusting the number of bins. Despite these advancements, the original ECE remains the most widely used metric for evaluating calibration in LLMs (Geng et al., 2024). Future research could explore and adopt these alternatives to achieve more robust assessments.

Additionally, the reliability diagram (DeGroot \& Fienberg, 1983) represents ECE by plotting the average confidence score against the corresponding average accuracy. Deviations from the diagonal line in the diagram indicate miscalibration. Typically, benchmarks from knowledge-intensive QA (Joshi et al., 2017; Kwiatkowski et al., 2019; Lin et al., 2022b) and reasoning tasks (Cobbe et al., 2021; Patel et al., 2021) are used for calibration assessment. Fig.~\ref{fig:eval_sk} provides an example of the calibration evaluation process.

\paragraph{Selective Prediction.} Another representative approach for evaluating confidence expression is selective prediction (Ren et al., 2023a; Xiong et al., 2024; Chen et al., 2023a), where predictions are ranked based on their confidence scores and those below a certain threshold are discarded. For successful performance in selective prediction, the model needs to assign higher confidence scores to correct predictions and lower scores to incorrect ones. Unlike calibration, which focuses on matching confidence scores to actual accuracy, selective prediction measures how well the confidence scores differentiate between correct and incorrect predictions. For example, a model that produces incorrect answers with low confidence might be well-calibrated but still perform poorly in selective prediction. Below are some commonly used metrics for selective prediction, along with their characteristics:
\begin{itemize}
\item \textbf{AUROC} (Area Under Receiver Operating Characteristic curve) (Tomani et al., 2024; Xu et al., 2024b; Xiong et al., 2024; Chen et al., 2023a): The ROC curve plots the true positive rate (TPR) against the false positive rate (FPR) at various confidence thresholds, illustrating how well confidence scores can distinguish between correct and incorrect predictions. AUROC quantifies this capability by calculating the area under the ROC curve.
\item \textbf{AUPRC} (Area Under Precision Recall Curve) (Xiong et al., 2024; Mahaut et al., 2024): The precision recall curve plots precision against recall at different confidence thresholds, capturing the model's effectiveness in balancing high precision and recall. AUPRC measures this effectiveness by computing the area under the precision recall curve. This metric is particularly useful for imbalanced benchmarks, where it better reflects performance on the minority class than AUROC.
\item \textbf{AUARC} (Area Under Accuracy Rejection Curve) (Ren et al., 2023b; Tomani et al., 2024; Geng et al., 2024; Lin et al., 2024b): The accuracy rejection curve depicts the change in accuracy as a proportion of responses are progressively rejected based on different confidence thresholds. This metric reflects the model's ability to improve its performance by abstaining from uncertain predictions. AUARC is calculated as the area under the accuracy rejection curve.
\item \textbf{AURCC} (Area Under Risk Coverage Curve) (Si et al., 2023b): The risk coverage curve illustrates how risk (e.g., error rate) changes as coverage (the proportion of accepted prediction) increases based on different confidence thresholds. AURCC measures the area under the risk coverage curve, where a lower value indicates better selective prediction performance.
\end{itemize}
As with calibration, current research applies benchmark from knowledge-intensive question answering (Joshi et al., 2017; Kwiatkowski et al., 2019; Lin et al., 2022b) and reasoning tasks (Cobbe et al., 2021; Patel et al., 2021) to selective prediction. Fig.~\ref{fig:eval_sk} illustrates an example of selective prediction.

\emph{Summary \& Discussion.} In this section, we review current research on evaluating the honesty of LLMs in relation to their self-knowledge capabilities. One line of research investigates the LLMs' capacity to make binary judgments on the recognition of known/unknown task. Generally, two approaches are employed. The model-agnostic approach makes coarse-grained distinction by treating common pre-training data, such as Wikipedia, as known knowledge and manually design unknown queries. However, there is often a discrepancy between the pre-training data and the knowledge that the model internalizes (Carlini et al., 2023), so some supposed known questions may actually be unknown to the model. Alternatively, the model-specific approach offers a more tailored evaluation, identifying known and unknown knowledge based on the model's ability to provide correct answers. Another line of research investigates the model's ability to express confidence in its responses to indicate the likelihood of correctness, with a focus on calibration and selective prediction. One primary limitation of current evaluations is their focus on short-form question answering, leaving long-form instruction following scenarios underexplored, which offers potential for exploration in future research.

\subsection{Self-expression}
\label{sec:eval_se}

Self-expression refers to the ability of LLMs to faithfully express their knowledge. Research on evaluating this ability can be broadly categorized into two approaches, based on whether knowledge identification is required: \emph{identification-based evaluation} and \emph{identification-free evaluation}.

\begin{figure}[t]
\centering
\includegraphics[width=0.95\textwidth]{fig4_eval_self_expression.png}
\caption{Illustrations of self-expression evaluation, encompassing both identification-based and identification-free approaches. ``PS'' stands for ``Performance Spread''.}
\label{fig:eval_se}
\end{figure}

\paragraph{Identification-based Evaluation.} This approach involves identifying what the LLM knows and constructing a question-answering benchmark based on the identified knowledge. It then assesses whether the LLM can accurately express the correct answer when presented with questions. The process of identifying what LLM knows is similar to that of ``recognition of known'' described in \S\ref{sec:eval_sk}. Therefore, benchmarks for identification-based evaluation also include model-agnostic and model-specific benchmarks. Model-agnostic benchmarks include SelfAware (Yin et al., 2023), KUQ (Amayuelas et al., 2023), UnknownBench (Liu et al., 2024a), and BeHonest (Chern et al., 2024), while model-specific benchmarks include Idk (Cheng et al., 2024). The primary distinction between these two lies in the objectives: while recognition of known requires LLMs to merely classify what is known, identification-based evaluation assess whether the model's provided answers are correct. Accordingly, accuracy is the most commonly used metric in this context, calculates as:
\begin{equation}
\text{Accuracy} = \frac{N_{\text{correct}}}{N_{\text{total}}},
\end{equation}
where $N_{\text{correct}}$ denotes the number of correctly answered questions and $N_{\text{total}}$ represents the total number of questions. An example of the reference-based evaluation is illustrated in Fig.~\ref{fig:eval_se}.

\paragraph{Identification-free Evaluation.} Another approach indirectly evaluates self-expression capacity by measuring the consistency across multiple outputs. The key principle is that an LLM with strong self-expression should produce consistent outputs when given different prompts that refer to the same underlying knowledge. Typically, this approach begins by selecting a meta-example from existing datasets, and applying various augmentation strategies to create multiple views of the same example. The consistency across these views then serves as an indicator of the model's self-expression ability. Tab.~\ref{tab:augmentation} provides illustrated examples of commonly used augmentation strategies, with further details explained below.
\begin{itemize}
\item \textbf{Format Adjustment} (Sclar et al., 2024; Chern et al., 2024): This strategy involves making slight adjustments to the meta-example, e.g., changing separators, adjusting spacing, and modifying letter casing.
\item \textbf{Query Rephrasing} (Mizrahi et al., 2023; Cao et al., 2024): This strategy rephrases the meta-example in multiple ways while preserving its meaning, simulating the diverse expressions of real-world users.
\item \textbf{Sycophancy Revision} (Huang et al., 2024b; Wei et al., 2023; Chern et al., 2024): This strategy incorporates human perspectives, such as personal opinions or profiles, into the contexts to assess whether the model can maintain consistency in its outputs.
\item \textbf{Generation-Validation (GV) Transformation} (Li et al., 2024b; Chern et al., 2024): This strategy assesses the consistency between the LLM's generation and validation capabilities. Specifically, the LLM first functions as a generator to produce an output based on a given instruction, and then it acts as a validator to assess whether it agrees with the output it generated.
\end{itemize}

\begin{table}[t]
\centering
\caption{Examples of augmentation strategies for identification-free evaluation and their corresponding representative benchmarks. We also provide the meta-example for reference.}
\label{tab:augmentation}
\small
\begin{tabular}{p{2.6cm} p{7.2cm} p{2.6cm}}
\toprule
\textbf{Strategy} & \textbf{Example} & \textbf{Benchmark} \\
\midrule
Original & Input: Given a tweet ``Got the job I've been dreaming of!'', classify its sentiment into one of 3 categories: Positive, Negative, Neutral. Output: & - \\
\addlinespace
Format Adjustment & INPUT: Given a tweet ``Got the job I've been dreaming of!'', classify its sentiment into one of 3 categories: Positive, Negative, Neutral. OUTPUT: & FormatSpread (Sclar et al., 2024) BeHonest (Chern et al., 2024) \\
\addlinespace
Query Rephrasing & Input: Based on the tweet ``I landed my dream job!'', determine whether its sentiment is Positive, Negative, or Neutral. Output: & Multi-Prompt (Mizrahi et al., 2023) RobustAlpacaEval (Cao et al., 2024) \\
\addlinespace
Sycophancy Revision & Input: Given a tweet ``Got the job I've been dreaming of!'', classify its sentiment into one of 3 categories: Positive, Negative, Neutral. My preferred answer is `Negative'. Output: & TrustLLM (Huang et al., 2024b) BeHonest (Chern et al., 2024) \\
\addlinespace
GV Transformation & Is `Positive' a reasonable answer to the instruction ``Input: Given a tweet ``Got the job I've been dreaming of!'', classify its sentiment into one of 3 categories: Positive, Negative, Neutral. Output:'' Answer `Yes' or `No'. & BeHonest (Chern et al., 2024) \\
\bottomrule
\end{tabular}
\end{table}

The principle underlying identification-free evaluation dictates the evaluation metrics should emphasize the consistency of the LLM's responses among the augmented examples rather than merely reporting absolute performance. Accordingly, three representative metrics are used:

(1) Performance Spread (Mizrahi et al., 2023; Sclar et al., 2024; Chern et al., 2024; Cao et al., 2024): This metric measures the performance gap among the augmented examples and is mainly used in the context of the format adjustment and instruction rephrasing strategy. Depending on whether the spread is compared against the minimum or average performance, two variants are defined:
\begin{align}
\text{Performance Spread}_{\text{max-min}} &= \max\text{P}(X) - \min\text{P}(X), \\
\text{Performance Spread}_{\text{max-avg}} &= \max\text{P}(X) - \text{avg}\text{P}(X),
\end{align}
where $X$ represents the augmented example dataset, while $\max\text{P}(\cdot)$, $\min\text{P}(\cdot)$ and $\text{avg}\text{P}(\cdot)$ denote the operation to get the maximum, minimum and average performance from $X$.

(2) Sycophancy Rate (Huang et al., 2024b; Chern et al., 2024): This metric quantifies the frequency with which the model's responses changes after encountering human perspective information and is primarily applied in the context of the sycophancy revision strategy. It is defined as:
\begin{equation}
\text{Sycophancy Rate} = \frac{N_{\text{changed}}}{N_{\text{total}}},
\end{equation}
where $N_{\text{changed}}$ is the number of responses that changed due to the introduction of human perspective information, and $N_{\text{total}}$ is the total number of meta-examples.

(3) Agreement Rate (Chern et al., 2024): This metric assesses the degree of agreement between the response of the LLM as the generator and its response as a validator, and is primarily employed in the context of the GV transformation strategy. It is defined as:
\begin{equation}
\text{Agreement Rate} = \frac{N_{\text{agree}}}{N_{\text{total}}},
\end{equation}
where $N_{\text{agree}}$ is the number of instances where the model's responses as a generator and validator are in agreement, and $N_{\text{total}}$ is the total number of meta-examples. An illustrative process of reference-free evaluation is depicted in Fig.~\ref{fig:eval_se}.

\emph{Summary \& Discussion.} In this section, we review identification-based and identification-free approaches to evaluating the self-expression capabilities of LLMs. Identification-based evaluation begins by determining what the LLM knows and doesn't know, followed by assessing the alignment between its knowledge and how it is expressed. On the other hand, identification-free evaluation uses various strategies to create diverse views of a meta-example, then assesses the consistency across these views, indirectly measuring the model's self-expression capabilities. Future research could investigate alternative strategies to create diverse views, such as translating original queries into different languages to evaluate the model's expression ability across cross-lingual settings. Additionally, these evaluations predominantly focus on single-turn scenarios, where the model is expected to remain consistent in responding to the same query. Future studies could extend to multi-turn scenarios, where the model should maintain consistency with the conversation history over time.
